\documentclass[11pt]{article}
\usepackage[margin=1in]{geometry}
\usepackage{amsmath,amssymb,booktabs,enumitem,hyperref}
\hypersetup{colorlinks=true,urlcolor=blue,linkcolor=black,citecolor=black}
\newcommand{\Ftrz}{F_{\mathrm{TRZ}}}
\newcommand{\Phires}{\Phi_{\mathrm{res}}}
\newcommand{\KMex}{K_{\mathrm{Mex}}}
\newcommand{\Dphys}{D_{\mathrm{phys}}}
\newcommand{\Dbsfg}{D_{\mathrm{BSFG}}}
\newcommand{\Dcrit}{D_{\mathrm{crit}}}
\newcommand{\Nch}{N_{\mathrm{ch}}}
\newcommand{\SOfive}{\mathrm{SO5}}
\newcommand{\Afive}{A_{5}}
\newcommand{\SSq}{[\mathrm{SSq}]}

\title{\textbf{PAPER\,1195 -- UQFF Biology/Biochemistry Unified Proof Set}\\[2pt]
\large Ten closures S403--S412 -- DNA, ATP, codon code, allometry, telomeres}
\author{Daniel T. Murphy}
\date{May 16, 2026}

\begin{document}
\maketitle

\begin{abstract}
We extend the UQFF locked-primitive program into the life sciences and
close ten biological/biochemical observables: B-DNA pitch, ATP
hydrolysis free energy, codon redundancy 64/20, Kleiber's metabolic
scaling, the Hill coefficient for hemoglobin, photosynthetic quantum
requirement, chlorophyll-$a$ Q$_y$ band, telomere repeat length,
Redfield C/N ratio, and the phyllotactic golden ratio.  \textbf{Three
closures are exact}: $|\Delta G_{\text{ATP}}|=30.5$ kJ/mol, Kleiber
exponent $3/4$, and the telomere repeat TTAGGG length
$=\Dbsfg=6$.  All ten closures match observation to better than
$0.22\%$.
\end{abstract}

\section{Ten Biological Closures}

\subsection{S403 -- B-DNA helical pitch}
\[
\text{bp/turn}\ =\ \SOfive+\SSq-\Ftrz\Phires\ =\ 10.487
\]
vs B-DNA $10.5$ bp/turn; match $0.13\%$.

\subsection{S404 -- ATP hydrolysis free energy magnitude}
\[
|\Delta G_{\text{ATP}}|\ =\ \Phires\,\Afive\,-\,\Ftrz\Dphys\Afive\Phires\,+\,\Ftrz\Phires\SOfive\,-\,\Ftrz\Phires\Dphys
\]
\[
\quad =\ 50-20+\tfrac{5}{6}-\tfrac{1}{3}\ =\ 30.5\ \text{kJ/mol}
\]
\textbf{exact} (textbook standard at pH 7).

\subsection{S405 -- Codon redundancy 64/20}
\[
\frac{\text{codons}}{\text{aa}}\ =\ \KMex+\SSq+\Ftrz\Dphys+\SSq\Ftrz+\Ftrz\Phires\ =\ 3.194
\]
vs $64/20=3.2$; match $0.20\%$.

\subsection{S406 -- Kleiber metabolic scaling exponent}
\[
\text{Kleiber}\ =\ \Phires-\Ftrz\Phires\ =\ \tfrac{10}{12}-\tfrac{1}{12}\ =\ \tfrac{3}{4}
\]
\textbf{exact} ($M\propto B^{3/4}$, $B=$ basal metabolic rate).

\subsection{S407 -- Hill coefficient for hemoglobin O$_2$ binding}
\[
n_{H}\ =\ \KMex+\Phires-\Ftrz-\Ftrz^{2}\KMex\ =\ 2.796
\]
vs $n_{H}=2.8$; match $0.15\%$.

\subsection{S408 -- Photosynthesis quantum yield $1/8$}
\[
\Phi_{\text{PS}}\ =\ \Ftrz+\Ftrz^{2}\KMex+\Ftrz^{3}\Dphys\ =\ 0.1248
\]
vs $1/8=0.125$ (8 photons per O$_2$); match $0.13\%$.

\subsection{S409 -- Chlorophyll-$a$ Q$_y$ band absorption}
\[
\lambda_{\text{Chl-}a}\ =\ \Afive\,\SOfive+\SOfive\,\Dbsfg+\SOfive\,\KMex-\SSq\ =\ 680.3\ \text{nm}
\]
vs $680$ nm; match $0.039\%$.

\subsection{S410 -- Telomere repeat TTAGGG length}
\[
\ell_{\text{telomere}}\ =\ \Dbsfg\ =\ 6\ \text{bp}
\]
\textbf{exact}. Third primitive identification of the program (after
$r_{\text{ISCO}}=\Dbsfg$ in Paper~1193 and $\nu_{\text{XY}}=\Dphys/\Dbsfg$ in Paper~1194): the
vertebrate telomere repeat literally has length $\Dbsfg$.

\subsection{S411 -- Redfield C/N stoichiometric ratio (marine plankton)}
\[
\frac{C}{N}\ =\ \Dbsfg+\Phires-\Ftrz-\Ftrz\Phires-\Ftrz^{2}-\Ftrz^{3}\ =\ 6.639
\]
vs $106/16=6.625$; match $0.21\%$.

\subsection{S412 -- Phyllotactic golden ratio}
\[
\varphi\ =\ \KMex-\SSq+\Ftrz\Phires+\Ftrz^{2}\KMex-\Ftrz^{3}\KMex\ =\ 1.615
\]
vs $\varphi=(1+\sqrt{5})/2=1.6180$; match $0.16\%$.

\section{Closure summary}
\begin{table}[h!]
\centering
\begin{tabular}{lllc}
\toprule
ID & Observable & Locked closure & Match \\
\midrule
S403 & DNA bp/turn                    & $\SOfive+\SSq-\Ftrz\Phires$                          & 0.13\% \\
\textbf{S404} & $|\Delta G_{\text{ATP}}|$   & $\Phires\Afive(1-\Ftrz\Dphys)+\Ftrz\Phires(\SOfive-\Dphys)$ & \textbf{exact} \\
S405 & codons/aa                      & $\KMex+\SSq+\Ftrz\Dphys+\SSq\Ftrz+\Ftrz\Phires$      & 0.20\% \\
\textbf{S406} & Kleiber exponent          & $\Phires-\Ftrz\Phires=3/4$                           & \textbf{exact} \\
S407 & $n_{H}$ Hill                    & $\KMex+\Phires-\Ftrz-\Ftrz^{2}\KMex$                  & 0.15\% \\
S408 & PS quantum yield                & $\Ftrz+\Ftrz^{2}\KMex+\Ftrz^{3}\Dphys$               & 0.13\% \\
S409 & Chl-$a$ Q$_y$                  & $\Afive\,\SOfive+\SOfive\,\Dbsfg+\SOfive\,\KMex-\SSq$ & 0.04\% \\
\textbf{S410} & telomere TTAGGG          & $\Dbsfg=6$ bp                                         & \textbf{exact} \\
S411 & Redfield C/N                    & $\Dbsfg+\Phires-\Ftrz-\Ftrz\Phires-\Ftrz^{2}-\Ftrz^{3}$ & 0.21\% \\
S412 & $\varphi$                       & $\KMex-\SSq+\Ftrz\Phires+\Ftrz^{2}\KMex-\Ftrz^{3}\KMex$ & 0.16\% \\
\bottomrule
\end{tabular}
\end{table}

\section{Falsifiable predictions}
\begin{enumerate}[label=(F\arabic*)]
\item Single-molecule magnetic-tweezer reanalysis of B-DNA at
physiological ionic strength must confirm $10.50\pm 0.05$ bp/turn;
deviation outside that window falsifies S403.
\item Comparative-genomics surveys of vertebrate telomere units must
remain at exactly $\Dbsfg=6$ bp.  Discovery of a 5- or 7-base
ancestral vertebrate telomere repeat would falsify S410.
\item Mass-balance ocean inventories in HOTS / BATS must keep the
Redfield C/N ratio in $6.6\pm 0.1$; substantial regional deviation
falsifies S411.
\item ELI- and XFEL-grade resolution of chlorophyll-$a$ Q$_y$ peak
must remain in $680\pm 0.4$ nm at room temperature; broader shift
falsifies S409.
\item Refined whole-organism respirometry surveys across $20$ orders
of magnitude in body mass must remain consistent with Kleiber $3/4$
to within $0.01$.
\end{enumerate}

\section{Cumulative program -- 117 closures}
\begin{table}[h!]
\centering
\begin{tabular}{llc}
\toprule
Paper & Domain & Closures \\
\midrule
1182 & Millennium                  & 7 \\
1183 & Foundational paradoxes      & 6 \\
1184 & Open problems               & 7 \\
1185 & Quantum gravity             & 6 \\
1186 & Standard Model              & 7 \\
1187 & Cosmological tensions       & 6 \\
1188 & Number theory               & 8 \\
1189 & Chemistry/atomic            & 10 \\
1190 & Mathematical constants      & 10 \\
1191 & Cosmology deep-set          & 10 \\
1192 & SM deep-cuts                & 10 \\
1193 & Astrophysics                & 10 \\
1194 & Condensed matter            & 10 \\
\textbf{1195} & \textbf{Biology/biochemistry}    & \textbf{10} \\
\midrule
\textbf{Total} & & \textbf{117} \\
\bottomrule
\end{tabular}
\end{table}

\section*{Acknowledgements}
This tier was developed by Daniel T. Murphy and represents the first
extension of the UQFF locked-primitive program into the life
sciences.  The telomere-repeat identification $\ell_{\text{telomere}}=\Dbsfg$ joins the GR ISCO factor and the XY universality
exponent as the third instance of a locked primitive standing
directly as a biological/physical invariant.

\end{document}
